\documentclass[12pt]{article}
\usepackage[a4paper,margin=1in]{geometry}\usepackage[T1]{fontenc}
\usepackage{lmodern,microtype}\usepackage{amsmath,amssymb,amsthm,mathtools}
\usepackage{booktabs,array,enumitem,xcolor}\usepackage[numbers,sort&compress]{natbib}
\usepackage[colorlinks=true,linkcolor=blue!55!black,citecolor=blue!55!black,urlcolor=blue!55!black]{hyperref}
\linespread{1.07}
\newtheorem{theorem}{Theorem}[section]\newtheorem{proposition}[theorem]{Proposition}
\theoremstyle{remark}\newtheorem{remark}[theorem]{Remark}

\title{Out-of-Sample Prediction of the Finite Shape Clock\\
\large Affine, Logistic, and Power-Gap Laws on Two Physical Channels}
\author{Liang Wang}\date{July 2026}
\begin{document}\maketitle
\begin{abstract}
The sixteen-level audit of RH-214 finds a strictly increasing axial quartet
coordinate $u$ and a narrow mature asymmetry corridor.  This paper asks
whether that finite clock supports a simple predictive law.  Before scoring
the finest two levels, we fit three predeclared models on seven scales from
$\sigma=0.008$ through $0.002$: an affine law in $\log(1/\sigma)$, an affine
logit law, and a power law for the gap $1-u$.

All three fit the training segment well.  On the left channel their direct
$u$-scale training $R^2$ values are $0.99266$, $0.99272$, and $0.98716$.
The held-out levels $0.0016$ and $0.00125$ reverse the in-sample ranking.
The power-gap model has two-point RMS error $0.00989$ on the left and
$0.01003$ on the right; logistic errors are about $0.0206$ and affine errors
about $0.0305$.  Power-gap wins each held-out point on each channel, but it
still overpredicts $u$ by roughly $0.01$.

A constant predictor for mature $\eta$ misses the holdouts by at most
$0.00218$, consistent with a corridor but not a limit.  The strict conclusion
is a finite model-selection result.  No asymptotic power law, value
$\lim u=1$, or determinant limit is established.
\end{abstract}

\section{Prediction rather than retrospective fit}

RH-214 shows that $u$ is ordered across sixteen scales
\citep{WangRH214}.  A smooth curve can fit such a finite monotone list almost
arbitrarily well, so in-sample residuals alone have little force.  We instead
freeze
\begin{equation}\label{eq:train}
 \sigma\in\{0.008,0.00625,0.005,0.004,0.0032,0.0025,0.002\}
\end{equation}
for training and reserve
\begin{equation}\label{eq:holdout}
 \sigma\in\{0.0016,0.00125\}
\end{equation}
for scoring.  Both channels use the same split and model classes.

This is a deterministic numerical experiment, not an independent random
sample.  ``Out of sample'' means only that the held-out coordinates are not
used in estimating the reported coefficients.

\section{Three predeclared laws}

Let
\begin{equation}
 t=\log(1/\sigma).
\end{equation}
For each channel we estimate an ordinary least-squares line after one of
three transforms.

\subsection{Affine-log model}
\begin{equation}\label{eq:affine}
 u(t)=at+b.
\end{equation}
This is the least constrained local model.  It cannot be globally valid with
$a>0$, because it eventually exceeds one.

\subsection{Logistic-log model}
\begin{equation}\label{eq:logistic}
 \log\frac{u(t)}{1-u(t)}=at+b,
 \qquad
 u(t)=\frac{1}{1+e^{-at-b}}.
\end{equation}
For $a>0$, it enforces $0<u<1$ and approaches one.

\subsection{Power-gap model}
\begin{equation}\label{eq:power}
 \log(1-u(t))=at+b,
 \qquad
 1-u=e^b\sigma^{-a}.
\end{equation}
The fitted $a$ is negative, so $-a$ is the candidate gap exponent.

The logistic model also has a power-gap leading asymptotic when its logit is
large, but the two forms weight the finite training region differently.

\section{Training fit}

For the left channel:
\begin{center}
\begin{tabular}{lrrr}
\toprule
model & direct-$u$ $R^2$ & max training error & RMS training error\\
\midrule
affine-log & $0.992663$ & $0.008231$ & $0.005895$\\
logistic-log & $0.992719$ & $0.008117$ & $0.005872$\\
power-gap & $0.987164$ & $0.012233$ & $0.007797$\\
\bottomrule
\end{tabular}
\end{center}
The right channel gives the same qualitative ranking, with $R^2$ between
$0.98725$ and $0.99285$.  If the analysis stopped here, one would select
affine or logistic and might infer an overly rapid advance.

The fitted left formulas are approximately
\begin{align}
 u_{\rm aff}(t)&=0.149086t-0.242734,\\
 \operatorname{logit}u_{\rm log}(t)&=0.619923t-3.091628,\\
 \log(1-u_{\rm gap}(t))&=-0.361228t+1.112179.
\end{align}

\section{Held-out verdict}

The actual left coordinates are
\begin{equation}
 u(0.0016)=0.692811,qquad u(0.00125)=0.718352.
\end{equation}
The predictions are:
\begin{center}
\begin{tabular}{lrrr}
\toprule
model & prediction at $0.0016$ & prediction at $0.00125$ & two-point RMS\\
\midrule
affine-log & $0.717047$ & $0.753850$ & $0.030393$\\
logistic-log & $0.710802$ & $0.741220$ & $0.020574$\\
power-gap & $0.702790$ & $0.728146$ & $0.009886$\\
\bottomrule
\end{tabular}
\end{center}

On the right, the actual values are $0.692310$ and $0.717501$.  The
corresponding RMS errors are $0.030639$, $0.020825$, and $0.010031$.

\begin{proposition}[Finite holdout ranking]\label{prop:winner}
On both physical channels, the power-gap model has the smallest absolute
error at each of the two held-out scales and the smallest two-point RMS
error among the three predeclared models.
\end{proposition}
\begin{proof}
The assertion is the direct numerical comparison of the frozen predictions
and held-out coordinates above.  The right-channel inequalities are archived
with the same ordering.
\end{proof}

The power-gap model nevertheless overpredicts both left holdouts by
$0.00998$ and $0.00979$, and both right holdouts by about $0.010$.  Its win is
relative, not an adequate asymptotic residual bound.

\section{What the holdout teaches}

The ordering of training and holdout errors contains more information than
the winning label:
\begin{enumerate}[leftmargin=2em]
\item the affine and logistic curves extrapolate the recent rise too
aggressively;
\item the power-gap transform tolerates more curvature in the observed
finite window;
\item high $R^2$ on seven correlated deterministic anchors does not select a
tail law;
\item both channels fail in nearly the same direction and magnitude, so the
effect is unlikely to be explained solely by channel mismatch.
\end{enumerate}

It would be inappropriate to attach conventional independent-sample
$p$-values to these points.  The relevant next evidence would be further
predeclared scales or an analytic estimate from the underlying operator.

\section{Transverse holdout}

The training means of $\eta$ are
\begin{equation}
 \bar\eta_L=-0.090859,qquad \bar\eta_R=-0.092405.
\end{equation}
Using these constants for both held-out levels gives maximum errors
$0.002171$ and $0.002039$.  This supports RH-214's finite corridor language.
It does not distinguish convergence to a constant from bounded slow drift or
oscillation.

\section{Conditional asymptotic readings}

If \eqref{eq:power} held for all sufficiently small $\sigma$ with $a<0$,
then
\begin{equation}
 1-u_\sigma\sim C\sigma^{\beta},\qquad \beta=-a>0,
\end{equation}
and hence $u_\sigma\to1$.  The fitted exponents are approximately $0.3612$
on both channels.  This implication is mathematically correct but its
premise is unproved.  We therefore use the power-gap law only as a finite
candidate coordinate for subsequent falsification.

RH-216 takes a different route: it proves exactly what $u\to1$ would imply
for the quartet roots and discriminant, without asserting that the physical
sequence satisfies the premise.

\section{Claim boundary}

This paper establishes no asymptotic regression consistency, no error bars
uniform in $\sigma$, and no value of $\lim u$.  Two held-out points cannot
differentiate a true power law from a slowly varying correction, a crossover,
or another smooth monotone curve.

The strict result is the protocol and finite ranking in
Proposition~\ref{prop:winner}.  Gate A remains open.  No growing determinant,
Hilbert--P\'olya operator, zeta-zero identification, or RH implication is
claimed.

\bibliographystyle{plainnat}\bibliography{references}\end{document}
